이 Part는 개별 구현을 관통하는 설계 철학을 정리하고, 이 철학이 \emph{실시간}보다
\emph{저장된 raw 데이터의 오프라인 재처리}에서 더 큰 효과를 낼 수 있음을 논한다.

\chapter{센서 상보성}

\section{레이더가 주는 것과 못 주는 것}
AWR1642는 \textbf{2 TX $\times$ 4 RX}의 1차원 가상 배열(8소자 ULA)이다. 따라서
\begin{itemize}
  \item \textbf{준다}: 정확한 거리, 방위각(azimuth), 도플러 속도, 반사 강도/SNR.
  \item \textbf{못 준다}: \textbf{고도각(elevation)}. 안테나가 수평으로만 배열되어
        있어 객체의 ``높이''를 분해하지 못한다. 또한 약반사 객체에서 검출이
        불안정(노이즈/허위검출)하다.
\end{itemize}
즉 레이더 검출은 본질적으로 \emph{수평면(거리--방위각) 위의 점}이며, 3D 공간의
완전한 위치를 스스로 확정하지 못한다.

\section{카메라(DA3)가 주는 것과 못 주는 것}
\begin{itemize}
  \item \textbf{준다}: 조밀한 픽셀별 깊이. 이미지의 세로축이 곧 높이이므로,
        깊이맵은 \emph{elevation을 암묵적으로 포함}한다. 또한 의미(클래스/마스크)와
        공간적으로 매끄러운(노이즈가 적은) 구조를 제공한다.
  \item \textbf{못 준다}: 절대 미터 스케일. 단안 깊이는 상대 관계만 신뢰할 수 있다.
\end{itemize}

\begin{designbox}{두 센서는 정확히 서로의 결손을 메운다}
레이더는 \emph{절대 거리}를 주지만 \emph{높이와 매끄러움}이 없고, DA3는
\emph{높이와 매끄러운 구조}를 주지만 \emph{절대 거리}가 없다. 본 시스템의 융합은
이 상보성을 양방향으로 활용한다: (1) 레이더로 DA3의 스케일을 미터로 고정하고,
(2) DA3로 레이더의 elevation 결손과 노이즈를 보완한다.
\end{designbox}

\chapter{융합의 세 가지 역할}

\section{역할 1 — 스케일 고정 (레이더 $\to$ DA3)}
신뢰도 높은(SNR$\ge$임계) 레이더 검출의 거리를 기준으로 객체별 비율(ratio)을
학습해 DA3 깊이를 실제 미터로 환산한다. 이것이 1부에서 구현한 핵심 경로다.

\section{역할 2 — elevation 보완 (DA3 $\to$ 레이더)}
레이더 검출점은 (거리, 방위각)만 가지므로 화면상 \emph{세로 위치}가 모호하다.
DA3 깊이맵과 카메라 투영을 쓰면, 같은 거리/방위각을 갖는 화면 열에서 \emph{어느
높이}에 그 객체가 있는지를 카메라가 채워준다. 즉 레이더 단독으로는 불가능한
``elevation을 가진 3D 점''을 카메라 깊이가 부여한다. 이는 레이더--카메라
외부 캘리브를 수평(yaw)뿐 아니라 \emph{높이 축까지} 확장하는 토대가 된다.

\section{역할 3 — 노이즈 보정 (상호 검증)}
레이더의 약검출/허위검출은 DA3 깊이와의 \emph{일관성}으로 걸러낼 수 있다.
객체 마스크 영역의 DA3 깊이가 가리키는 기대 거리(expected)와 동떨어진 레이더
점은 게이트에서 배제되고, 이상치 비율은 스케일 학습에서 거부된다. 반대로
DA3의 스케일 드리프트는 신뢰도 높은 레이더 검출이 교정한다. 두 센서가 서로의
오류를 견제한다.

\chapter{왜 오프라인 재처리에서 더 효과적인가}

본 시스템은 실시간으로 동작하지만, 같은 융합 철학은 \textbf{저장된 raw ADC
(\texttt{.bin})를 사후에 다시 읽을 때 잠재력이 더 크다}. 시스템이 raw ADC와
타임스탬프, 웹캠 프레임을 모두 저장하도록 설계한 이유가 여기에 있다.

\section{실시간의 제약}
\begin{itemize}
  \item DA3 재추론이 7초 주기(ComfyUI 부하)라 매 프레임 최신 깊이를 못 쓴다.
  \item 라이브 루프는 0.3초 예산 안에서 동작해야 해 무거운 정밀화가 어렵다.
  \item 본질적으로 \emph{인과적}(과거만 사용)이라 미래 프레임을 활용한 평활이 불가.
\end{itemize}

\section{오프라인이 여는 것}
\begin{itemize}
  \item \textbf{전 프레임 고품질 깊이}: 시간 제약 없이 모든 프레임에 ESRGAN +
        대형 DA3를 적용해 조밀하고 정확한 깊이를 얻는다.
  \item \textbf{양방향(비인과적) 시간 필터}: 과거뿐 아니라 \emph{미래 프레임}까지
        보고 레이더 거리/속도를 평활할 수 있어, 실시간 EMA보다 훨씬 강하게
        노이즈를 제거한다(정지 물체의 거리 널뛰기를 거의 완전히 흡수).
  \item \textbf{사후 파라미터 튜닝}: CFAR 임계, 게이트, 비율 신뢰 기준을 데이터
        전체를 보고 최적화한 뒤 일괄 재처리.
  \item \textbf{elevation 완성}: 모든 레이더 검출점에 DA3 elevation을 부여해,
        레이더가 단독으로 만들 수 없는 \emph{완전한 3D 포인트 집합}(거리·방위각은
        레이더, 높이는 카메라)을 구성. 객체 단위 캘리브와 3D 배치가 정밀해진다.
  \item \textbf{누적 외부 캘리브}: 전체 시퀀스의 대응점을 한꺼번에 사용해 yaw/높이
        정합을 더 정확히 추정.
\end{itemize}

\begin{designbox}{핵심 주장}
실시간 시스템은 ``레이더로 카메라 깊이의 스케일을 고정''하는 데까지를 견고하게
달성했다. 그러나 \emph{elevation이 없는 레이더 보드를 카메라 깊이로 보완하고,
양방향 필터로 레이더 노이즈를 보정하며, 객체 단위 3D 캘리브를 정밀화}하는 일은,
시간 예산과 인과성 제약이 사라지는 \textbf{저장된 \texttt{.bin}의 오프라인
재처리}에서 비로소 온전히 실현된다. 본 시스템이 raw ADC와 타임스탬프를 끝까지
보존하도록 설계한 것은 이 오프라인 확장을 염두에 둔 결정이다.
\end{designbox}
